\chapter{الحساب}\label{ch:g12:arith}

يدرس الحساب \omterm{def:g10:numbers:sets}{الأعداد الصحيحة}: \omterm{def:g12:arith:divides}{قابلية القسمة}، والأعداد الأولية، والبواقي.
وقد عُدّ طويلًا أنقى الرياضيات البحتة، وهو اليوم يحمي كل دفعة
عبر الإنترنت: فنظام التعمية RSA يقوم على مبرهنات بيزو وغاوس و
فيرما المبرهَن عليها في هذا الفصل.

\section{قابلية القسمة والقسمة الإقليدية}

\begin{definition}[قابلية القسمة]\label{def:g12:arith:divides}
ليكن $a, b \in \Z$. نقول إن $b$ \emph{يقسم} $a$\index{قابلية القسمة}، ويُكتب
$b \mid a$، إذا وُجد $k \in \Z$ حيث $a = kb$. ونقول أيضًا إن $a$
\emph{مضاعف} للعدد $b$.
\end{definition}

\begin{proposition}\label{prop:g12:arith:divprops}
إذا كان $c \mid a$ و $c \mid b$، فإن $c$ \omterm{def:g12:arith:divides}{يقسم} كل تركيب صحيح
$au + bv$ (حيث $u, v \in \Z$). وإذا كان $a \mid b$ و $b \mid a$ مع
$a,b \in \N$، فإن $a = b$. وإذا كان $a \mid b$ و $b \neq 0$، فإن $\abs a \leq \abs b$.
\end{proposition}

\begin{proof}
اكتب $a = kc$ و $b = lc$: فيكون $au + bv = (ku + lv)c$. وتنتج النقطتان الأخريان
من $\abs{a} = \abs{k}\,\abs{b}$ مع $\abs k \geq 1$ عندما يكون
$b = ka \neq 0$.
\end{proof}

\begin{theorem}[القسمة الإقليدية]\label{thm:g12:arith:euclid}
\index{قسمة إقليدية}
ليكن $a \in \Z$ و $b \in \N^*$. يوجد زوج وحيد
$(q, r) \in \Z \times \N$ بحيث
\[
a = bq + r \qquad\text{و}\qquad 0 \leq r < b .
\]
ويكون $q$ هو \emph{خارج القسمة} و $r$ هو \emph{الباقي}.
\end{theorem}

\begin{proof}
\emph{الوجود.} لمجموعة مضاعفات $b$ التي لا تتجاوز $a$ عنصر
أكبري $bq$ (فهي غير خالية \omterm{def:g12:seq:bounded}{ومحدودة من الأعلى})؛ ضع $r = a - bq$.
وبالأكبرية، $b(q+1) > a$، إذن $0 \leq r < b$.
\emph{الوحدانية.} إذا كان $bq + r = bq' + r'$ مع $0 \leq r, r' < b$، فإن
$b(q - q') = r' - r$ و $\abs{r' - r} < b$: ومضاعف للعدد $b$ قيمته المطلقة
أصغر من $b$ يجب أن يكون $0$، إذن $r = r'$ و $q = q'$.
\end{proof}

\section{الموافقات}

\begin{definition}[الموافقة]\label{def:g12:arith:congruence}
ليكن $n \in \N^*$. يكون عددان صحيحان $a, b$ \emph{متوافقين بترديد
$n$}\index{موافقة}، ويُكتب $a \equiv b \pmod n$، إذا كان $n \mid (a - b)$
--- وبصيغة مكافئة، إذا كان للعددين $a$ و $b$ الباقي نفسه في \omterm{thm:g12:arith:euclid}{القسمة
الإقليدية} على $n$.
\end{definition}

\begin{proposition}[التوافق مع العمليات]\label{prop:g12:arith:congops}
إذا كان $a \equiv b \pmod n$ و $c \equiv d \pmod n$، فإن
\[
a + c \equiv b + d, \qquad
ac \equiv bd, \qquad
a^k \equiv b^k \ (k \in \N) \pmod n .
\]
\end{proposition}

\begin{proof}
\omterm{def:g12:arith:divides}{يقسم} $n$ المقدار $(a-b) + (c-d) = (a+c) - (b+d)$، و
$ac - bd = a(c - d) + d(a - b)$ مضاعف للعدد $n$ أيضًا. وتنتج قاعدة القوة
بالتراجع من قاعدة الجداء.
\end{proof}

\begin{method}[حساب القوى بترديد $n$]\label{met:g12:arith:powers}
لحساب $a^k \bmod n$، ردّ \omterm{def:g11:seq:arithmetic}{الأساس} بترديد $n$، ثم ابحث عن قوة صغيرة
للعدد $a$ توافق $\pm1$، واستعملها لطيّ الأس. فمثلًا
$2^{100} \bmod 7$: بما أن $2^3 = 8 \equiv 1 \pmod 7$ و
$100 = 3\times33 + 1$، فإن
\[
2^{100} = \left(2^{3}\right)^{33} \times 2 \equiv 1^{33}\times 2 = 2 \pmod 7 .
\]
\end{method}

\section{القاسم المشترك الأكبر وبيزو وغاوس}

\begin{definition}[القاسم المشترك الأكبر]\label{def:g12:arith:gcd}
ليكن $a, b$ عددين صحيحين غير معدومين معًا. \emph{القاسم المشترك
الأكبر}\index{قاسم مشترك أكبر} $\gcd(a, b)$ هو أكبر \omterm{def:g10:numbers:sets}{عدد صحيح} \omterm{def:g12:arith:divides}{يقسم} $a$
و $b$ معًا. وعندما يكون $\gcd(a,b) = 1$، يقال إن $a$ و $b$
\emph{أوليان فيما بينهما}\index{أوليان فيما بينهما}.
\end{definition}

\begin{proposition}[خوارزمية إقليدس]\label{prop:g12:arith:euclidalgo}
\index{خوارزمية إقليدس}
إذا كان $a = bq + r$ (حيث $b \neq 0$)، فإن $\gcd(a, b) = \gcd(b, r)$. ومنه فإن تكرار
\omterm{thm:g12:arith:euclid}{القسمة الإقليدية} يحسب $\gcd(a,b)$: \omterm{def:g12:arith:gcd}{فالقاسم المشترك الأكبر} هو آخر باقٍ
غير معدوم.
\end{proposition}

\begin{proof}
كل قاسم مشترك للعددين $a$ و $b$ \omterm{def:g12:arith:divides}{يقسم} $r = a - bq$
(\cref{prop:g12:arith:divprops})، ومنه فهو قاسم مشترك للعددين $b$ و $r$؛
وبالعكس، لأن $a = bq + r$. فللزوجين القواسم المشتركة نفسها،
إذن \omterm{def:g12:arith:gcd}{القاسم المشترك الأكبر} نفسه. وتنتهي الخوارزمية لأن البواقي
تكوّن \omterm{def:g12:seq:sequence}{متتالية} \omterm{def:g10:functions:variations}{متناقصة} تمامًا من \omterm{def:g10:numbers:sets}{الأعداد الصحيحة} غير السالبة.
\end{proof}

\begin{example}\label{ex:g12:arith:euclidalgo}
$\gcd(252, 198)$: $252 = 198 + 54$؛ و $198 = 3\times54 + 36$؛
و $54 = 36 + 18$؛ و $36 = 2 \times 18 + 0$. ومنه $\gcd(252,198) = 18$.
\end{example}

\begin{theorem}[مساواة بيزو]\label{thm:g12:arith:bezout}
\index{مساواة بيزو}
ليكن $a, b$ عددين صحيحين غير معدومين معًا، وليكن $d = \gcd(a,b)$. يوجد
$u, v \in \Z$ بحيث
\[
au + bv = d .
\]
وخاصة، يكون $a$ و $b$ أوليين فيما بينهما إذا وفقط إذا كان $au + bv = 1$ من أجل
عددين صحيحين $u, v$ ما.
\end{theorem}

\begin{proof}
شغّل \omterm{prop:g12:arith:euclidalgo}{خوارزمية إقليدس} بالمقلوب: فكل باقٍ تركيب صحيح
للباقيين السابقين، والمعطيان الابتدائيان $a, b$ تركيبان من
نفسيهما؛ وبالتعويض النازل، يكون آخر باقٍ غير معدوم $d$
تركيبًا صحيحًا للعددين $a$ و $b$. (وفي \cref{ex:g12:arith:euclidalgo}:
$18 = 54 - 36 = 54 - (198 - 3\times54) = 4\times54 - 198 =
4(252 - 198) - 198 = 4\times252 - 5\times198$.)

وأما التكافؤ: فإذا كان $\gcd(a,b) = 1$، أعطت مساواة بيزو $u, v$؛ وبالعكس،
كل قاسم مشترك للعددين $a$ و $b$ \omterm{def:g12:arith:divides}{يقسم} $au + bv = 1$، فيفرض
$\gcd(a,b) = 1$.
\end{proof}

\begin{theorem}[مبرهنة غاوس المساعدة]\label{thm:g12:arith:gauss}
\index{مبرهنة غاوس المساعدة}
ليكن $a, b, c \in \Z$. إذا كان $a \mid bc$ و $\gcd(a, b) = 1$، فإن $a \mid c$.
\end{theorem}

\begin{proof}
تعطي مساواة بيزو $au + bv = 1$؛ فاضرب في $c$: $acu + bcv = c$. وحدّا
الطرف الأيسر مضاعفان للعدد $a$ (والثاني لأن $a \mid bc$)، ومنه
كذلك $c$.
\end{proof}

\begin{corollary}\label{cor:g12:arith:coprimeprod}
إذا كان $a \mid c$ و $b \mid c$ و $\gcd(a,b) = 1$، فإن $ab \mid c$.
\end{corollary}

\begin{proof}
اكتب $c = ak$. ومن $b \mid ak$ و $\gcd(a,b)=1$، تعطي مبرهنة غاوس المساعدة
أن $b \mid k$، وليكن $k = bl$؛ عندئذٍ $c = abl$.
\end{proof}

\section{الأعداد الأولية}

\begin{definition}[العدد الأولي]\label{def:g12:arith:prime}
يكون \omterm{def:g10:numbers:sets}{العدد الصحيح} $p \geq 2$ \emph{أوليًا}\index{عدد أولي} إذا كانت قواسمه
الموجبة الوحيدة هي $1$ و $p$.
\end{definition}

\begin{proposition}\label{prop:g12:arith:primedivides}
لكل \omterm{def:g10:numbers:sets}{عدد صحيح} $n \geq 2$ قاسم أولي؛ وإذا لم يكن $n$ \omterm{def:g12:arith:prime}{أوليًا}، فله
قاسم أولي $\leq \sqrt n$. وإذا قسم \omterm{def:g12:arith:prime}{عدد أولي} $p$ جداءً $ab$، فإن
$p \mid a$ أو $p \mid b$ (وهي \emph{مبرهنة إقليدس المساعدة}).
\end{proposition}

\begin{proof}
أصغر قاسم $d \geq 2$ للعدد $n$ أولي (فأي قاسم فعلي للعدد $d$
سيكون قاسمًا أصغر للعدد $n$). وإذا كان $n = de$ مركّبًا مع
$2 \leq d \leq e$، فإن $d^2 \leq de = n$، إذن $d \leq \sqrt n$. وأما مبرهنة إقليدس
المساعدة: فإذا كان $p \nmid a$، كان $\gcd(p, a) = 1$ (لأن قواسم $p$ الوحيدة هي
$1$ و $p$)، فتعطي مبرهنة غاوس المساعدة أن $p \mid b$.
\end{proof}

\begin{theorem}[إقليدس]\label{thm:g12:arith:infiniteprimes}
الأعداد الأولية لا نهائية العدد.
\end{theorem}

\begin{proof}
إذا أُعطيت أي قائمة منتهية $p_1, \dots, p_k$ من الأعداد الأولية، فاعتبر
$N = p_1 p_2 \cdots p_k + 1$. عندئذٍ \omterm{def:g12:arith:divides}{يقسم} \omterm{def:g12:arith:prime}{عدد أولي} $p$ ما العدد $N$؛ لكن لا \omterm{def:g12:arith:divides}{يقسم} أي $p_i$
العدد $N$ (فالباقي $1$)، إذن $p$ \omterm{def:g12:arith:prime}{عدد أولي} ليس في القائمة. فلا
قائمة منتهية تستنفد الأعداد الأولية.
\end{proof}

\begin{theorem}[المبرهنة الأساسية في الحساب]\label{thm:g12:arith:fta}
\index{المبرهنة الأساسية في الحساب}
كل \omterm{def:g10:numbers:sets}{عدد صحيح} $n \geq 2$ جداء أعداد أولية، وهذا التحليل
وحيد إلى غاية \omterm{def:g10:coordgeom:system}{ترتيب} العوامل:
\[
n = p_1^{\alpha_1} p_2^{\alpha_2} \cdots p_r^{\alpha_r},
\qquad p_1 < p_2 < \dots < p_r \text{ أعداد أولية},\ \alpha_i \geq 1 .
\]
\end{theorem}

\begin{proof}
\emph{الوجود}، بالتراجع القوي: فالعدد $n$ الأولي تحليله نفسه؛
وإلا فإن $n = de$ مع $2 \leq d, e < n$، ويتحلّل كلاهما بفرضية
التراجع. و\emph{الوحدانية}: لنفترض
$p_1\cdots p_s = q_1 \cdots q_t$ (أعدادًا أولية، مع السماح بالتكرار). فحسب
مبرهنة إقليدس المساعدة، \omterm{def:g12:arith:divides}{يقسم} $p_1$ عددًا ما $q_j$، وبما أنه أولي، فإن $p_1 = q_j$؛
فاختصر وكرّر. وينطبق التحليلان حدًا بحد.
\end{proof}

\begin{theorem}[مبرهنة فيرما الصغرى]\label{thm:g12:arith:fermat}
\index{مبرهنة فيرما الصغرى}
ليكن $p$ عددًا \omterm{def:g12:arith:prime}{أوليًا} وليكن $a \in \Z$ حيث $p \nmid a$. عندئذٍ
\[
a^{p-1} \equiv 1 \pmod p .
\]
ومن أجل كل $a \in \Z$ (بلا فرض الأولية فيما بينهما)، $a^p \equiv a \pmod p$.
\end{theorem}

\begin{proof}
اعتبر \omterm{def:g10:numbers:sets}{الأعداد الصحيحة} $p - 1$ التالية $a, 2a, 3a, \dots, (p-1)a$ بترديد $p$. فلا واحد منها
$\equiv 0$ (إذ لو كان $p \mid ka$ مع $1 \leq k \leq p-1$، لفرضت مبرهنة إقليدس المساعدة
أن $p \mid k$، وهذا مستحيل)، وهي متمايزة مثنى مثنى بترديد $p$ (إذ لو كان
$ka \equiv la$، لكان $p \mid (k - l)a$، إذن $p \mid k - l$، إذن $k = l$).
ومنه فهي، بترديد $p$، الأعداد $1, 2, \dots, p-1$ \omterm{def:g10:coordgeom:system}{بترتيب} ما.
وبضرب كل الموافقات:
\[
a^{p-1}\,(p-1)! \equiv (p-1)! \pmod p .
\]
وبما أن $p$ لا \omterm{def:g12:arith:divides}{يقسم} أيًا من $1, \dots, p-1$، فإن الاستعمال المتكرر لمبرهنة إقليدس المساعدة
يسمح باختصار $(p-1)!$، فيبقى $a^{p-1} \equiv 1$. وتنتج الصورة الثانية
بالضرب في $a$ (وهي بديهية عندما يكون $p \mid a$).
\end{proof}

\begin{example}[تطبيق في التعمية]\label{ex:g12:arith:rsa}
تجعل مبرهنة فيرما رفع القوى بترديد $n$ قابلًا للعكس عندما تُختار
الأسس اختيارًا مناسبًا --- وهذا قلب نظام التعمية \emph{RSA}.
فمع $p, q$ عددين أوليين كبيرين و $n = pq$، يُنشَر $n$ وأس
$e$؛ ويكون التعمية $x \mapsto x^e \bmod n$. أما فكّ التعمية فيقتضي أسًا
$d$ حيث $ed \equiv 1 \pmod{(p-1)(q-1)}$، ولا يستطيع حسابه إلا من يعرف $p$
و $q$ --- واستعادة $p, q$ من $n$ تعني تحليل عدد
طوله مئات الأرقام، وهو ما لا تفعله أي خوارزمية معروفة في زمن
معقول.
\end{example}

\section{تمارين}

\begin{exercise}[$\star$]\label{exo:g12:arith:1}
احسب خارج قسمة وباقي \omterm{thm:g12:arith:euclid}{القسمة الإقليدية} للعدد $2026$ على
$17$، وللعدد $-2026$ على $17$.
\end{exercise}

\begin{exercise}[$\star$]\label{exo:g12:arith:2}
ما باقي $7^{100}$ بترديد $10$؟ (أي ما آخر رقم في
$7^{100}$؟)
\end{exercise}

\begin{exercise}[$\star$]\label{exo:g12:arith:3}
باستعمال \omterm{prop:g12:arith:euclidalgo}{خوارزمية إقليدس}، احسب $\gcd(1071, 462)$، وأوجد عددين صحيحين
$u, v$ حيث $1071u + 462v = \gcd(1071, 462)$.
\end{exercise}

\begin{exercise}[$\star$]\label{exo:g12:arith:4}
بيّن أنه من أجل كل $n \in \Z$، يكون $n^2$ موافقًا $0$ أو $1$ بترديد
$4$. واستنتج أن عددًا صحيحًا $\equiv 3 \pmod 4$ لا يكون أبدًا مجموع مربعين.
\end{exercise}

\begin{exercise}[$\star\star$]\label{exo:g12:arith:5}
بيّن أنه من أجل كل $n \in \N$، يقبل $n(n+1)(2n+1)$ القسمة على $6$.
\end{exercise}

\begin{exercise}[$\star\star$]\label{exo:g12:arith:6}
حل في $\Z$ \omterm{def:g12:arith:congruence}{الموافقة} $5x \equiv 3 \pmod{11}$.
(إرشاد: أوجد مقلوب $5$ بترديد $11$.)
\end{exercise}

\begin{exercise}[$\star\star$]\label{exo:g12:arith:7}
حل في $\Z \times \Z$ \omterm{def:g10:algebra:equation}{المعادلة} الديوفانتية
\[
17x - 40y = 1,
\]
ثم صِف كل حلول $17x - 40y = 6$.
\end{exercise}

\begin{exercise}[$\star\star$]\label{exo:g12:arith:8}
بيّن أن $\sqrt2$ \omterm{ex:g10:numbers:classify}{عدد أصم}، باستعمال وحدانية \omterm{def:g10:algebra:expand}{التحليل إلى عوامل}
أولية (قارن أس $2$ في طرفي $a^2 = 2b^2$).
\end{exercise}

\begin{exercise}[$\star\star\star$]\label{exo:g12:arith:9}
ليكن $p$ عددًا \omterm{def:g12:arith:prime}{أوليًا}.
\begin{enumerate}
  \item بيّن أنه من أجل $1 \leq k \leq p - 1$، \omterm{def:g12:arith:divides}{يقسم} $p$ المقدار $\dbinom{p}{k}$.
        (إرشاد: استعمل $k\binom pk = p\binom{p-1}{k-1}$ و \cref{exo:g12:comb:7}،
        ومبرهنة غاوس المساعدة.)
  \item استنتج، بالتراجع على $a \geq 0$، برهانًا آخر لمبرهنة فيرما
        الصغرى على الصورة $a^p \equiv a \pmod p$.
\end{enumerate}
\end{exercise}

\begin{exercise}[$\star\star\star$]\label{exo:g12:arith:10}
\emph{(مسألة البواقي الصينية.)} أوجد كل \omterm{def:g10:numbers:sets}{الأعداد الصحيحة} $n$ التي تحقق
\[
n \equiv 2 \pmod 3, \qquad n \equiv 3 \pmod 5, \qquad n \equiv 2 \pmod 7 .
\]
(إرشاد: حل الشرطين الأولين، ثم أدخل الثالث؛ ومعاملات بيزو
تساعد.)
\end{exercise}

\section{مسألة: الشيفرات السرية وأرقام التحقق}

\begin{problem}[{مسألة نهاية الأسبوع --- الموافقات تحرس كل
رمز شريطي وكل بطاقة ائتمان، ومبرهنة فيرما الصغرى تدير
قفل أسرار العالم}]
\label{pb:g12:arith:1}
تباهى هاردي سنة 1940 بأن نظرية الأعداد ``غير ملوَّثة''
بالتطبيقات. وبعد ثمانين سنة، يكذّبه كل صفير رمز شريطي، وكل
دفعة ببطاقة ائتمان، وكل رسالة معمّاة
--- وبأدوات هذا الفصل بالضبط: الموافقات
(\cref{prop:g12:arith:congops})، ومقلوبات بيزو
(\cref{thm:g12:arith:bezout})، ومبرهنة فيرما الصغرى
(\cref{exo:g12:arith:9}). وتتحقق هذه المسألة من الشيفرات، وتكسر
نسخة لعبة من القفل، وتتعلم لماذا يصمد القفل الحقيقي.

\medskip
\noindent\textbf{الجزء الأول --- التمكّن من الموافقات.}
\begin{enumerate}
  \item احسب $2026 \bmod 7$؛ ثم آخر رقم في
        $7^{100}$ (أوجد دورة قوى $7$ بترديد
        $10$).
  \item رفع القوى السريع (\cref{met:g12:arith:powers}):
        احسب $5^{117} \bmod 13$ (انطلق من
        $5^2 \equiv -1$).
  \item حل $3x \equiv 5 \pmod 7$.
  \item شغّل \omterm{prop:g12:arith:euclidalgo}{خوارزمية إقليدس} على $(97, 35)$، ثم عوّض بالمقلوب
        لإيجاد عددين صحيحين $u, v$ حيث $97u + 35v = 1$، ثم
        استنتج مقلوب $35$ بترديد $97$.
  \item صُغ بدقة متى يكون $a$ قابلًا للقلب بترديد $n$، و
        أي مبرهنة تسلّم المقلوب.
\end{enumerate}

\medskip
\noindent\textbf{الجزء الثاني --- أرقام التحقق.}
\begin{enumerate}[resume]
  \item الترميز ISBN ذو العشرة أرقام: يجب أن تحقق الأرقام العشرة $d_1 \dots d_{10}$ لشيفرة
        كتاب العلاقة
        $10d_1 + 9d_2 + \dots + 2d_9 + 1d_{10} \equiv 0
        \pmod{11}$. تحقق من الشيفرة الحقيقية
        $0\,306\,40615\,2$.
  \item برهن على أن مخطط ISBN يكشف \emph{كل}
        خطأ في رقم واحد: فإذا تغيّر رقم بمقدار
        $d \not\equiv 0$، تغيّر المجموع المرجَّح بالمقدار $w d$
        حيث $1 \leq w \leq 10$ --- فلماذا لا يمكن أن يكون هذا أبدًا
        $\equiv 0 \pmod{11}$
        (\cref{thm:g12:arith:gauss})؟
  \item برهن على أنه يكشف أيضًا كل تبادل بين رقمين
        متجاورين (متمايزين). ثم فسّر سرّ التصميم: أي
        خاصية للعدد $11$ جعلت البرهانين يفلحان،
        وما الذي قد يسوء مع الترديد $10$؟
  \item ترجّح الرموز الشريطية EAN ذات الثلاثة عشر رقمًا الأرقام بالأوزان $1, 3, 1, 3, \dots$
        بترديد $10$. احسب رقم التحقق الذي يكمل
        $978\,2940199\,05$. وأي التبادلات المتجاورة \emph{يعجز}
        هذا الترميز عن كشفها؟ (ومتى يكون
        $2(a - b) \equiv 0 \pmod{10}$؟)
  \item تستعمل بطاقات الائتمان مخطط لوهن: من اليمين، ضاعف
        كل رقم ثانٍ (مع طرح $9$ عندما يتجاوز الضعف $9$)،
        ثم اجمع كل شيء، واشترط أن يكون المجموع مضاعفًا للعدد
        $10$. تحقق من الرقم الاختباري
        $4539\,1488\,0343\,6467$.
  \item في جملة واحدة: ماذا اشترى الترديد الأولي للترميز ISBN
        مما لا يستطيعه EAN ومخطط لوهن المقيَّدان بالعدد $10$؟
\end{enumerate}

\medskip
\noindent\textbf{الجزء الثالث --- قفل فيرما.}
\begin{enumerate}[resume]
  \item مصيدة قبل الكنز: احسب
        $2^{10} \bmod 341$، واستنتج $2^{340} \bmod 341$ --- ثم
        حلّل $341$. فماذا يقول هذا المثال (وهو
        \emph{عدد فيرما الأولي الكاذب}) عن استعمال مبرهنة فيرما
        الصغرى اختبارًا للأولية؟
  \item نظام RSA مصغَّرًا: خذ $p = 3$ و $q = 11$، فيكون
        $n = 33$ و $(p-1)(q-1) = 20$؛ والأس العلني هو
        $e = 3$. أوجد الأس السري $d$ حيث
        $3d \equiv 1 \pmod{20}$ (بمنهج السؤال 4).
  \item عمِّ الرسالة $m = 4$: احسب
        $c = m^3 \bmod 33$.
  \item فكّ التعمية: احسب $c^d \bmod 33$ (باستعمال
        $c \equiv -2 \pmod{33}$) واستعد الرسالة.
  \item لماذا ينجح فكّ التعمية دائمًا: بيّن أن
        $m^{21} \equiv m$ بترديد $3$ وبترديد $11$ معًا
        (بمبرهنة فيرما الصغرى في كل عالم)، ثم اختم
        بترديد $33$ (إذ تلصق \cref{thm:g12:arith:gauss} الموافقتين).
        وأين دخلت الصورة الخاصة $1 + 20k$ للمقدار $21 = ed$؟
  \item أمان القفل: الجميع يعرف $n$ و $e$؛ و
        استعادة $d$ تقتضي $(p-1)(q-1)$، ومنه عوامل
        $n$. وعددنا $33$ يتحلّل بالنظر --- فلماذا يحمي المخطط
        نفسه، مع $n$ من ستمئة رقم، مصارف
        العالم؟ (جملة واحدة عن اللاتناظر بين
        الضرب والتحليل.)
\end{enumerate}

\medskip
\noindent\textbf{الجزء الرابع --- الكلاسيكيات.}
\begin{enumerate}[resume]
  \item عدّ الجنود الصيني القديم (قارن
        \cref{exo:g12:arith:10}): عدد من الجنود يترك
        الباقي $2$ عند اصطفافهم صفوفًا من $3$، والباقي $3$ عند
        اصطفافهم صفوفًا من $5$. أوجد كل الأعداد الممكنة، وفسّر
        لماذا يكون الجواب وحيدًا بترديد $15$.
  \item براهين من سطر واحد أخيرًا: انطلاقًا من $10 \equiv 1 \pmod 9$،
        برهن على أن كل عدد يوافق مجموع أرقامه
        بترديد $9$؛ ومن $10 \equiv -1 \pmod{11}$، استخرج
        قاعدة المجموع المتناوب من أجل $11$. (وقد برهن الكتاب
        السابق على هاتين بجبر صريح --- فاعجب
        من الضغط.)
  \item الخاتمة --- هاردي في مواجهة الرمز الشريطي: لخّص
        عدّة الفصل (حساب الموافقات، ومقلوبات
        بيزو، ومبرهنة فيرما الصغرى، ولصق الترديدين الأوليين
        فيما بينهما) وأين انطبق كل منها في هذه
        المسألة؛ ثم أعطِ الحكم الحديث على
        ``غير ملوَّثة''.
\end{enumerate}
\end{problem}
